\section{Boundary Stability Simulation}
\label{sec:boundary}

\subsection{Simulation Methodology}

For each state with a projected seat change, we:
\begin{enumerate}
\item Run \ar{} with the 2020 seat count $k_{20}$ on the 2020 census
      tract adjacency graph with 2020 population weights.
\item Run \ar{} with the 2030 projected seat count $k_{30}$ on the same
      2020 graph (holding geography fixed to isolate the seat-change effect;
      a full 2030 simulation would require the 2030 census, not yet available).
\item Compute the Hamming distance between the two plans (after canonicalization).
\item Identify the districts with the largest boundary changes.
\end{enumerate}

Step 2 uses the 2020 graph because the 2030 census data are not yet available.
This means we are measuring the boundary disruption from the seat count change
alone, not from the combination of seat change and population redistribution.
The actual 2030 disruption will be larger (population redistribution adds
further disruption).

\subsection{Results}

\begin{table}[h]
\centering
\caption{Boundary stability simulation results. Hamming distance
         $d_{\rm Ham}$ between 2020-seat and 2030-seat plans on the
         2020 graph. ``Most affected'' is the district with the largest
         boundary change.}
\label{tab:boundary-stability}
\begin{tabular}{lcccl}
\toprule
State & $k_{20}$ & $k_{30}$ & $d_{\rm Ham}$ & Change type \\
\midrule
Texas       & 38 & 41 & 0.23 & Tree restructure (composite $\to$ prime) \\
California  & 52 & 50 & 0.14 & Tree restructure (different composite) \\
Florida     & 28 & 29 & 0.19 & Tree restructure (composite $\to$ prime) \\
New York    & 26 & 25 & 0.12 & Tree restructure (different composite) \\
Arizona     & 9  & 10 & 0.09 & Minor tree change \\
Georgia     & 14 & 15 & 0.08 & Tree restructure (different composite) \\
N. Carolina & 14 & 15 & 0.08 & Tree restructure (different composite) \\
Illinois    & 17 & 16 & 0.15 & Prime $\to$ four-level tree \\
Michigan    & 13 & 12 & 0.11 & Prime $\to$ three-level tree \\
Pennsylvania& 17 & 16 & 0.14 & Prime $\to$ four-level tree \\
\bottomrule
\end{tabular}
\end{table}

\subsection{Interpretation}

The Hamming distances range from 0.08 (Georgia: minor structural change)
to 0.23 (Texas: dramatic prime transition).
For context:
\begin{itemize}
\item A single \recom{} step changes $\approx 1/k$ of tracts.
      For Texas ($k = 38$): $d_{\rm Ham} \approx 0.026$ per step.
\item A single seed change within the current tree structure (B.7)
      changes $d_{\rm Ham} \approx 0.01$--$0.05$.
\end{itemize}

The Texas reapportionment ($d_{\rm Ham} = 0.23$) is equivalent to
approximately $0.23 / 0.026 \approx 9$ \recom{} steps.
This is a large disruption relative to algorithmic mixing, but small
relative to politically-motivated redistricting, which typically changes
30--60\% of tracts.

\subsection{Which Districts Are Most Affected?}

For Texas, the transition from a 2-level tree (two halves of 19) to a
flat 41-way tree most affects:
\begin{itemize}
\item The boundary between the two halves of the 2-level tree (the
      ``spine'' of the current 38-district map).
\item Districts near the state's centroid (where the two subtrees meet).
\item Urban districts in Houston and Dallas metro areas (where the
      prime-19 flat split is replaced by the global 41-way partition).
\end{itemize}

Rural districts are less affected because their geography constrains
both the 38- and 41-district maps similarly.

\paragraph{Geographic breakdown of tract changes (Texas).}
Of the tracts that change district assignment when Texas gains 3 seats
($38 \to 41$): 43\% are in urban cores (Houston, Dallas, Austin), 31\%
suburban, 26\% rural---consistent with the population growth patterns
driving the seat gain.

\subsection{Seed Variance Around Reported Disruption}
\label{sec:seed-variance}

The Hamming distances in Table~\ref{tab:boundary-stability} are measured
at seed $= 0$ (the DIA default seed).
Across $T = 600$ seeds, median disruption is $\pm 0.8$ percentage points
from the reported values.
The qualitative ordering of states by disruption magnitude is stable
across all seeds.
